\scanpage{194}
2) On peut aussi regarder le \unsure{hensélisé|hensilisé}
$\widetilde S$ de $S$, et les hom
\[
\mathrm{Pic}(S-a)\longrightarrow\mathrm{Pic}(\widetilde S-\widetilde a)-\mathrm{Pic}(\widehat S-\widehat a).
\]
Il est possible que la dernière flèche
écrite soit toujours un isomorphisme.
On le voit immédiatement, en tous cas,
lorsque $X$ est de dim $1$ (\unsure{i.e.|càd} $S$ de
dim $2$), par le raisonnement habituel
sur l'éclaté \ldots\,. \unsure{Donc|Dans} l'exemple précédent
des courbes planes donne autant
d'exemples où \unsure{les|le} \unsure{gpes|groupes|groupe} de Picard local de
$S$ et celui de son hensélisé ne sont
pas isomorphes \ldots\,.
% archive sheet number: 193
